% !TeX program = xelatex
% Complete transcription in scan order: Adobe Scan Oct 21, 2026 (1)_1.jpg--_28.jpg.
% The source wording and numbering are retained, including repeated numbers.
\ifdefined\settextfont
  \def\IPStandalone{0}
\else
  \def\IPStandalone{1}
  \documentclass[11pt,a4paper]{article}
  \usepackage[margin=1.7cm]{geometry}
  \usepackage{amsmath,amssymb,enumitem}
  \usepackage{xepersian}
  \IfFontExistsTF{Amiri}{\settextfont{Amiri}}{\settextfont{FreeSerif}}
  \setlatintextfont{FreeSerif}
  \setlength{\parindent}{0pt}
\fi
\providecommand{\en}[1]{\lr{#1}}
\newcommand{\IPpage}[1]{\par\clearpage\noindent{\footnotesize\en{Note book\quad Subject 27\quad Year:\ \_\_\quad Month:\ \_\_\quad Date:\ \_\_}\par}\subsection*{صفحهٔ #1}}
\newcommand{\IPfoot}[1]{\par\medskip{\footnotesize\en{BEHROO DAFTER 181}\hfill #1\par}}
\newenvironment{IPpoints}{\begin{itemize}[leftmargin=*,itemsep=.25em]}{\end{itemize}}
\newenvironment{IPenglish}{\begin{latin}\begin{itemize}[leftmargin=*,itemsep=.25em]}{\end{itemize}\end{latin}}
\ifnum\IPStandalone=1 \begin{document}\fi
% pic01.jpg
\IPpage{۱}
\section*{\en{Intracranial Pathology}}
\begin{enumerate}[label=\arabic*.,leftmargin=*]
\item در مغز \en{MRI} تصویربرداری انتخابی در اکثر موارد هست.
\item آناتومی نرمال مغز در \en{CT}
\begin{IPpoints}
\item بطن چهارم: به صورت یک آب هیپودنس دیده می‌شود.
\item در خلف بطن چهارم نیمکره‌های مخچه قرار دارند.
\item در قدام بطن چهارم، \en{Pons} قرار دارد.
\end{IPpoints}
\item تنتوریوم سربلی دو فضای اینفراتنتوریال و سوپراتنتوریال را از هم جدا می‌کند.
\begin{IPpoints}\item بطن چهار، مخچه\end{IPpoints}
\item سیسترن اینترپدانکولر
\begin{IPpoints}\item در \en{Midbrain} قرار دارد.\item سربرال پدانکل‌ها را از هم جدا می‌کند.\end{IPpoints}
\item سیسترن سوپراسلار
\begin{IPpoints}\item ظاهر ۵ یا ۶ پر مثل ستاره\item قدام سیسترن اینترپدانکولر\end{IPpoints}
\item فیشرهای \en{sylvian}
\begin{IPpoints}\item حاوی \en{CSF}\item لوب تمپورال را از لوب‌های فرونتال و پریتال جدا می‌کند.\end{IPpoints}
\end{enumerate}
\IPfoot{۱}
% pic02.jpg
\IPpage{۲}
\begin{enumerate}[label=\arabic*.,leftmargin=*]\setcounter{enumi}{6}
\item هسته لنتی‌فرم: دو بخش
\begin{IPpoints}\item پوتامن در لترال\item گلوبوس پالیدوس در مدیال\end{IPpoints}
\item بطن سوم
\begin{IPpoints}\item \en{slit-like}\item در خط میدلاین\end{IPpoints}
\end{enumerate}
\noindent ۸. غده پینه‌آل
\begin{IPpoints}\item در خلف بطن سوم قرار دارد.\item در خلف آن \en{Quadrigeminal plate cistern}\end{IPpoints}
\noindent ۹. \en{Corpus Callosum}
\begin{IPpoints}\item دو لوب مغز را به هم متصل می‌کند.\item سقف بطن‌های جانبی را می‌سازد.\item انتهای قدامی آن: \en{genu}\item انتهای خلفی آن: \en{splenium}\end{IPpoints}
\noindent ۱۰. \en{Basal ganglia}
\begin{IPenglish}\item Subthalamic nucleus\item Substantia nigra\item Globus pallidus\item Putamen\item Caudate nucleus\item Putamen + Caudate nucleus: Striatum\end{IPenglish}
\noindent ۱۱. شاخ‌های قدامی بطن‌های جانبی سر هسته کودیت را در بر می‌گیرند.
\begin{IPpoints}\item این دو شاخ در میدلاین توسط \en{septum pellucidum} از هم جدا می‌شوند.\end{IPpoints}
\IPfoot{۲}
% pic03.jpg
\IPpage{۳}
\noindent ۱۲. شاخ‌های تمپورال بطن‌های جانبی
\begin{IPpoints}\item بسیار کوچک هستند.\item تحتانی‌تر هستند و در دل لوب تمپورال قرار دارند.\end{IPpoints}
\noindent ۱۳. شاخ‌های خلفی (اکسی‌پیتال) بطن‌های جانبی
\begin{IPpoints}\item در دل لوب اکسی‌پیتال قرار دارند.\end{IPpoints}
\noindent ۱۴. فوقانی‌ترین بخش سیستم بطنی: بادی بطن‌های جانبی

\noindent ۱۵. داس مغزی (\en{Falx cerebri})
\begin{IPpoints}\item درون فیشر \en{interhemispheric}\item معمولاً در بزرگسالان کلسیفیه است.\end{IPpoints}
\noindent ۱۶. کورتکس
\begin{IPpoints}\item سولکوس (\en{grooves})\item جایروس (\en{elevations})\end{IPpoints}
\noindent ۱۷. در یک \en{CT} بدون کنتراست
\begin{IPpoints}\item هر چیز سفید: \en{Bone}، \en{Blood}، \en{Metal}\end{IPpoints}
\noindent ۱۸. کلسیفیکیشن‌های غیرپاتولوژیک
\begin{IPenglish}\item Pineal gland\item Basal ganglia\item Choroid plexus\item Falx and Tentorium\end{IPenglish}
\IPfoot{۳}
% pic04.jpg
\IPpage{۴}
\noindent ۱۹. ساختارهایی که در \en{CT} بعد از تزریق کنتراست \en{enhance} می‌شوند:
\begin{IPenglish}\item Venous sinuses\item Choroid plexus\item Pituitary gland and stalk\end{IPenglish}
\noindent ۲۰. آرتیفکت
\begin{IPpoints}\item فلز\item \en{Dental fillings}\item کلیپس آنوریسم\item گلوله\end{IPpoints}
\noindent ۲۱. \en{MRI} نسبت به \en{CT}: حدوداً حساسیت بیشتری دارد.
\begin{IPpoints}\item برای شناسایی کلسیفیکاسیون درون ضایعه و ارزیابی \en{Cortical bone}: \en{less sensitive}\end{IPpoints}
\subsection*{تصویربرداری انتخابی}
\begin{enumerate}[label=\arabic*.,leftmargin=*]
\item \en{Acute Stroke}
\begin{IPpoints}\item \en{Choice: D-W MRI}\item \en{Other: Non contrast CT}: افتراق هموراژیک از ایسکمیک\end{IPpoints}
\item \en{Acute, Severe Headache}
\begin{IPpoints}\item \en{Choice: N.C-CT}: تشخیص \en{SAH}\item \en{Other: MRA or CTA}: اگر \en{SAH} تأیید شد برای پیدا کردن آنوریسم\end{IPpoints}
\end{enumerate}
\IPfoot{۴}
% pic05.jpg
\IPpage{۵}
\begin{latin}
\begin{enumerate}[label=\arabic*.,leftmargin=*]\setcounter{enumi}{2}
\item Chronic Headache
\begin{itemize}\item Choice: MRI with and without contrast\item Other: CT with and without contrast\end{itemize}
\item Seizures
\begin{itemize}\item Choice: MRI with and without contrast with Thin Section images through hippocampi (Hippocampus protocol)\item Patients with history of childhood seizures\item Other: CT with and without contrast\end{itemize}
\item Blood
\begin{itemize}\item Choice: Noncontrast CT\item Other: US for patients\end{itemize}
\item Head Trauma
\begin{itemize}\item Choice: NC-CT\item Other: MRI: Diffuse Axonal Injury\end{itemize}
\item Extracranial Carotid Disease
\begin{itemize}\item Choice: Doppler US\item Other: MRA, CTA (Preoperative Stenosis Evaluation)\end{itemize}
\item Hydrocephalus
\begin{itemize}\item Choice: MRI (Initial Study)\item Other: CT (Follow-up)\end{itemize}
\end{enumerate}
\end{latin}
\IPfoot{۵}
% pic06.jpg
\IPpage{۶}
\begin{latin}
\begin{enumerate}[label=\arabic*.,leftmargin=*]\setcounter{enumi}{8}
\item Vertigo and Dizziness
\begin{itemize}\item Choice: Contrast-enhanced MRI\item MRA or MRI with auditory canal protocol\end{itemize}
\item Masses
\begin{itemize}\item Choice: MRI with and without contrast\item Other: Contrast-enhanced CT\end{itemize}
\item Change in mental status
\begin{itemize}\item Choice: MRI with and without contrast\item Other: NC-CT\end{itemize}
\end{enumerate}
\end{latin}
\subsection*{دنسیته‌های \en{CT}}
\noindent \en{Hypodense}
\begin{IPpoints}\item \en{Fat}\item \en{Air (Sinuses)}\item \en{Water (CSF)}\item هماتوم مزمن ساب‌دورال\end{IPpoints}
\noindent \en{Isodense}
\begin{IPpoints}\item \en{Normal brain}\item هماتوم تحت حاد ساب‌دورال\end{IPpoints}
\noindent \en{Hyperdense}
\begin{IPenglish}\item Metal\item Iodine\item Calcium\item Blood and hemorrhage\end{IPenglish}
\IPfoot{۶}
% pic07.jpg
\IPpage{۷}
\noindent ۱. در اکثر بیماران با \en{Pacemaker} می‌توان از \en{MRI} استفاده کرد.
\subsection*{مغز در سکانس‌های \en{MRI}}
\begin{latin}
\begin{enumerate}[label=\arabic*.,leftmargin=*]
\item Bright on $T_1$: Fat, Gad, High pro, Subacute hem, Melanin
\item Dark on $T_1$: Air, Calc, Water, chronic hem, Acute hem
\item Bright on $T_2$: Water, Hyperacute hem, Late subacute hem
\item Dark on $T_2$: Air, Fat, Calc, high pro, Early subacute, Acute, chronic hem
\end{enumerate}
\end{latin}
\noindent ۱. در \en{CT} برای بررسی استخوان: \en{Bone Window}

\noindent ۲. تروماهای حاد به سر
\begin{IPpoints}
\item تصویربرداری انتخابی: \en{CT} بدون کنتراست
\item هدف این است که آیا ضایعه تهدیدکننده حیات قابل درمانی وجود دارد یا خیر.
\item بررسی سه نکته: \en{Blood}، \en{Mass effect}، \en{Displacement}، \en{Compression}
\item به کجا بنگریم؟ \en{Ventricles}، \en{Basal Cisterns}، \en{Sulci}
\end{IPpoints}
\IPfoot{۷}
% pic08.jpg
\IPpage{۸}
\noindent ۸. برای بررسی وجود خون: به کجا نگاه کنیم؟
\begin{IPenglish}\item Basal Cisterns\item Sylvian and interhemispheric Fissures\item Ventricles\item Subdural and epidural space\item Brain parenchyma\end{IPenglish}
\subsection*{شکستگی جمجمه}
\begin{enumerate}[label=\arabic*.,leftmargin=*]
\item انواع
\begin{IPenglish}\item Linear\item Depressed\item Basilar\end{IPenglish}
\item شکستگی \en{Linear}
\begin{IPpoints}\item شایع‌ترین فرم شکستگی\item اهمیت بالینی کم\item معمولاً در نواحی تمپورال و پریتال اتفاق می‌افتند.\end{IPpoints}
\item شکستگی \en{Depressed}
\begin{IPpoints}\item بیشتر در نواحی فرونتوتمپورال\item اگر بخش شکسته به زیر بخش سالم برود: جراحی\end{IPpoints}
\item شکستگی \en{Basilar}
\begin{IPpoints}\item مهم‌ترین شکستگی‌ها\item یک شکستگی خطی در قاعده جمجمه است.\item ممکن است دورا هم به طور همزمان پاره شده باشد: \en{CSF leak}\end{IPpoints}
\end{enumerate}
\IPfoot{۸}
% pic09.jpg
\IPpage{۹}
\noindent ۵. شک به شکستگی \en{Basilar}
\begin{IPenglish}\item Traumatic Pneumocephalus\item Fluid in the mastoid air cells\item Air fluid level in the sphenoid sinus\end{IPenglish}
\subsection*{شکستگی‌های صورت}
\begin{enumerate}[label=\arabic*.,leftmargin=*]
\item تصویربرداری انتخابی: \en{CT}
\item خط شکستگی باید به چند اسلایس متوالی دیده شود. با یک اسلایس نمی‌گوییم که شکستگی دیده شد.
\item \en{Blow-out Frac}
\begin{IPpoints}
\item شایع‌ترین شکستگی اربیتال
\item نیرو مستقیماً به اربیت وارد شده است.
\item انواع: شکستگی دیواره اینفریور: به داخل سینوس ماگزیلاری؛ شکستگی دیواره مدیال: به داخل سینوس اتموئید
\item ممکن است عضله رکتوس تحتانی در شکستگی گیر کند: محدودیت \en{Upward gaze}، \en{Diplopia}
\item ممکن است هوا در اربیت باشد: آمفیزم اربیتال
\item تشخیص: مایع (خون) در سینوس ماگزیلاری؛ گیر افتادن چربی و عضله؛ آمفیزم اربیتال؛ دیدن خط شکستگی
\end{IPpoints}
\end{enumerate}
\IPfoot{۹}
% pic10.jpg
\IPpage{۱۰}
\noindent ۴. \en{Tripod Frac}
\begin{IPpoints}
\item به دنبال ترومای بلانت به گونه
\item نسبتاً شایع است.
\item جدا شدن زایگوما: به دنبال جدا شدن سوچور فرونتوزایگوماتیک + شکستگی کف اربیت + شکستگی دیواره لترال سینوس ایپسی‌لترال ماگزیلاری
\item ممکن است آمفیزم اربیت دیده شود.
\end{IPpoints}
\subsection*{خونریزی داخل جمجمه (\en{ICH})}
\noindent ۱. انواع \en{ICH}
\begin{IPpoints}\item هماتوم اپی‌دورال\item هماتوم ساب‌دورال (\en{SDH})\item خونریزی اینتراسربرال\item خونریزی ساب‌آراکنوئید (\en{SAH})\end{IPpoints}
\noindent ۲. هماتوم اپی‌دورال
\begin{IPpoints}
\item اغلب به علت آسیب به شریان یا ورید میدل مننژیال
\item معمولاً ناشی از ترومای بلانت
\item در ۹۵\% موارد با شکستگی جمجمه همزمان هستند؛ عموماً استخوان تمپورال
\item ممکن است در پی جدا شدن \en{Dural Venous Sinuses} ایجاد شوند.
\end{IPpoints}
\IPfoot{۱۰}
% pic11.jpg
\IPpage{۱۱}
\noindent ۳. تشخیص هماتوم اپی‌دورال
\begin{IPenglish}\item high density, extra axial, biconvex lens-shaped mass\item most often in temporoparietal region\end{IPenglish}
\begin{IPpoints}\item از سوچور لاین‌ها عبور نمی‌کنند.\item برخلاف هماتوم ساب‌دورال می‌توانند از تنتوریوم عبور کنند.\end{IPpoints}
\noindent ۴. ممکن است یک هماتوم اسکالپ روی جمجمه دیده شود.

\noindent ۵. هماتوم ساب‌دورال (\en{SDH})
\begin{IPpoints}
\item از اپی‌دورال شایع‌تر است.
\item معمولاً با شکستگی همراهی ندارد.
\item معمولاً در پی آسیب‌های دسلریشن و یا به دنبال سقوط
\item معمولاً ناشی از آسیب به \en{Bridging Veins}
\begin{IPenglish}\item Cross from cerebral cortex to the venous sinus of the brain\end{IPenglish}
\item خونریزی در فضای بالقوه بین دوراماتر و آراکنوئید جمع می‌شود.
\end{IPpoints}
\noindent ۶. یک \en{SDH} حاد معمولاً نشان‌دهنده یک آسیب به پارانشیم مغز است.
\IPfoot{۱۱}
% pic12.jpg
\IPpage{۱۲}
\noindent ۷. \en{SDH}: احتمال افزایش اینتراکرانیال پرشر وجود دارد.

\noindent ۸. \en{Recognizing SDH (acute)}
\begin{IPenglish}
\item CT: Crescent shaped
\item Extracerebral
\item high attenuation
\item Concave inward
\item May cross suture lines
\item May enter the interhemispheric fissure
\item Do not cross mid line
\end{IPenglish}
\noindent ۸. \en{SDH} با گذر زمان و بعد از ترکیب با \en{CSF}: \en{Isoattenuating}، \en{Isodense}
\begin{IPenglish}\item Compressed or absent sulci or displaced sulci\end{IPenglish}
\noindent ۹. در \en{SDH} بعد از یک هفته و به دنبال ته‌نشین شدن سلول‌ها: \en{Fluid--Fluid level}

\noindent ۱۰. \en{Chronic SDH}
\begin{IPpoints}\item پس از ۳ هفته از آسیب\item \en{usually low density} (تیره‌تر از بافت مغز)\end{IPpoints}
\subsection*{هماتوم اینتراسربرال}
\noindent ۱. علل: تروما، آنوریسم، بیماری‌های آتروماتوز عروق کوچک، واسکولیت، مالفورمیشن عروقی (کاورنوما، فیستول شریانی--وریدی، مالفورمیشن شریانی--وریدی)، ترومبوز وریدی، آمیلوئید
\IPfoot{۱۲}
% pic13.jpg
\IPpage{۱۳}
\noindent ۲. انواع
\begin{IPpoints}\item \en{Coup}: در نقطه برخورد\item \en{Contrecoup}: در طرف مقابل\item عموماً به دنبال تروما\end{IPpoints}
\noindent ۳. \en{Coup Injuries}: عموماً به دنبال پاره شدن عروق کوچک داخل مغزی

\noindent ۴. \en{Contrecoup Inj}: عموماً به دنبال آسیب اسلریشن/دسلریشن

\noindent ۵. هر دو نوع می‌توانند کانتیوژن مغزی ایجاد کنند.

\noindent ۶. کانتیوژن هموراژیک
\begin{IPpoints}\item خونریزی + ادم\item معمولاً در \en{inf. frontal lobe} و \en{ant. temporal lobe} نزدیک سطح مغز\end{IPpoints}
\noindent ۷. تروما همچنین می‌تواند کانتیوژن غیرهموراژیک نیز ایجاد کند.

\noindent ۸. یافته‌های \en{CT} برای هماتوم اینتراسربرال در طول زمان تغییر می‌کند.
\begin{IPpoints}\item \en{MRI} گزینه مناسب‌تری است.\end{IPpoints}
\noindent ۹. \en{CT finding of traumatic intracerebral hemorrhage}
\begin{latin}
\begin{enumerate}[label=\alph*),leftmargin=*]
\item Hemorrhagic contusions: Multiple, small, well demarcated areas of high attenuation
\item Hypodense rim from Edema
\item Intraventricular blood
\end{enumerate}
\end{latin}
\IPfoot{۱۳}
% pic14.jpg
\IPpage{۱۴}
\begin{latin}
\begin{enumerate}[label=\alph*),leftmargin=*]\setcounter{enumi}{3}
\item Mass effect
\begin{itemize}
\item Compression of the ventricles
\item Displacement of the third ventricle and septum pellucidum
\item Can produce severe brain or vascular damage
\item With sufficient mass effect: Transtentorial and subfalcine brain herniation and death
\end{itemize}
\end{enumerate}
10. Brain Herniation
\begin{enumerate}[label=\alph*),leftmargin=*]
\item Extracranial Herniation
\item Sphenoid Herniation
\item Foramen magnum/tonsilar Herniation
\item Transtentorial Herniation
\item Subfalcine Herniation
\end{enumerate}
\end{latin}
\subsection*{آسیب آکسون منتشر (\en{DAI})}
\begin{enumerate}[label=\arabic*.,leftmargin=*]
\item علت کما طولانی به دنبال تروما سر
\item بدترین پیش‌آگهی در آسیب‌های سر
\item بیهوش شدن در لحظه آسیب
\end{enumerate}
\IPfoot{۱۴}
% pic15.jpg
\IPpage{۱۵}
\noindent ۴. در \en{DAI}، کورپوس کالوزوم عموماً درگیر است.

\noindent ۵. در \en{DAI}، \en{CT} اولیه ممکن است نرمال باشد.
\begin{IPpoints}\item \en{MRI} تصویربرداری انتخابی می‌باشد.\end{IPpoints}
\noindent ۶. \en{MRI findings in DAI}
\begin{IPenglish}
\item Small petechial hemorrhage may be bright on $T_1$
\item Multiple bright areas on $T_2$
\begin{itemize}\item Temporal or parietal cervicomedullary junction\item Corpus Callosum\end{itemize}
\end{IPenglish}
\noindent ۷. ابزار دقیق‌تر:
\begin{IPenglish}\item Gradient susceptibility MRI sequences\item Punctate dark areas: petechial hemorrhage\end{IPenglish}
\subsection*{افزایش فشار داخل جمجمه (\en{Increased ICP})}
\begin{enumerate}[label=\arabic*.,leftmargin=*]
\item علائم: ادم پاپیلا، سردرد، دیپلوپیا
\item علل
\begin{IPpoints}\item ادم مغزی: افزایش سایز مغز\item هیدروسفالی: افزایش سایز بطن‌ها\end{IPpoints}
\item ادم مغزی؛ علل
\begin{IPpoints}\item تروما\item توده\item \en{HTN}: به علت خونریزی و سکته\end{IPpoints}
\end{enumerate}
\IPfoot{۱۵}
% pic16.jpg
\IPpage{۱۶}
\noindent ۴. انواع ادم مغزی
\begin{IPpoints}
\item وازوژنیک
\begin{IPpoints}\item تجمع مایع در خارج از سلول‌ها\item بدخیمی و عفونت\item به علت نفوذپذیری غیرطبیعی \en{BBB}\item بیشتر \en{white matter} را تحت تأثیر قرار می‌دهد.\end{IPpoints}
\item سایتوتوکسیک
\begin{IPpoints}\item ادم سلولی\item مرتبط با ایسکمی مغز\item مرتبط با مرگ سلولی\item \en{white matter} و \en{gray matter} را درگیر می‌کند.\end{IPpoints}
\end{IPpoints}
\noindent ۵. \en{Recognizing cerebral edema}
\begin{IPpoints}
\item از بین رفتن تفاوت بین ماده خاکستری و سفید در سایتوتوکسیک
\item افاسمان سولکوس‌ها
\item کامپرشن بطن‌ها
\item هرنیشن مغز خصوصاً اگر افاسمان سیسترن بازال رخ داده باشد.
\end{IPpoints}
\subsection*{سکته}
\noindent ۱. \en{Acute loss of neurologic function}

\noindent ۲. تشخیص سکته بالینی است.
\IPfoot{۱۶}
% pic17.jpg
\IPpage{۱۷}
\noindent ۳. علل انجام تصویربرداری در سکته
\begin{IPpoints}
\item آیا علل دیگری برای اختلالات نورولوژیک وجود دارد؟ (مثل تومور)
\item آیا خون وجود دارد؟ (افتراق سکته هموراژیک از ایسکمیک)
\item برای تعیین انفارکت
\end{IPpoints}
\noindent ۴. اغلب سکته‌ها زمینه آمبولیک دارند.
\begin{IPpoints}\item با منشأ اینترنال کاروتید یا کامان کاروتید\item حتی می‌توانند منشأ قلبی و قوس آئورت داشته باشند.\end{IPpoints}
\noindent ۵. درصدی از سکته‌ها، ترومبوتیک هستند.
\begin{IPpoints}\item ترومبوز \en{MCA} شایع است.\item ترومبوز کاروتید، ورتبروبازیلار، شریان‌های داخل مغزی\end{IPpoints}
\noindent ۶. انواع سکته
\begin{IPpoints}\item ایسکمیک: شایع‌تر؛ درمان با \en{tPA}\item هموراژیک\end{IPpoints}
\noindent ۷. در \en{NC-CT} سکته هموراژیک بلافاصله و سکته ایسکمیک بعد از چند ساعت قابل تشخیص هستند.

\noindent ۸. \en{DW-MRI}، انفارکت ناشی از سکته را بعد از حدود \en{20--30 min} نشان می‌دهد.
\IPfoot{۱۷}
% pic18.jpg
\IPpage{۱۸}
\noindent ۹. سکته ایسکمیک: شایع‌ترین علت بیماری ترومبوآمبولیک به دنبال آترواسکلروز

\noindent ۱۰. \en{Vascular watershed areas}: حساس‌ترین نقاط به ایسکمی

\noindent ۱۱. شایع‌ترین یافته سکته غیرهموراژیک حاد در \en{CT}: نرمال بودن

\noindent ۱۲. \en{Ischemic stroke on CT}
\begin{latin}
\begin{enumerate}[label=\alph*),leftmargin=*]
\item 12 to 24 hours: Indistinct area of low attenuation in a vascular distribution.
\item $>24$ h: Lesion with mass effect that peaks at 3--5 days and disappears by 2--4 weeks.
\item 72 h: Contrast enhancement
\end{enumerate}
\end{latin}
اگرچه کنتراست معمولاً استفاده نمی‌شود.
\begin{latin}
\begin{enumerate}[label=\alph*),leftmargin=*]\setcounter{enumi}{3}
\item $>4$ w: Mass effect disappears, well-circumscribed low attenuation lesion with no contrast enhancement
\end{enumerate}
13. Vascular territories of stroke
\begin{itemize}
\item MCA: Common
\item Entire cortex surface, basal ganglia
\item Contralateral hemiparesis or hypesthesia
\item Ipsilateral hemianopsia, gaze preference toward the side of the lesion
\item agnosia is common, weakness of the arm and face
\item aphasia if the lesion in the dominant hemisphere
\end{itemize}
\end{latin}
\IPfoot{۱۸}
% pic19.jpg
\IPpage{۱۹}
\begin{latin}
\noindent Anterior Cerebral artery: Uncommon
\begin{itemize}
\item All of the frontal and parietal lobes
\item anterior 4/5 of the corpus callosum
\item frontobasal cerebral cortex
\item anterior diencephalon
\item Disinhibition with preservation of speech
\item Primitive reflexes (grasping or sucking)
\item Altered mental status
\item impaired judgment
\item contralateral weakness (greater in legs)
\end{itemize}
\noindent Posterior Cerebral artery
\begin{itemize}
\item portions of midbrain, subthalamic nucleus, basal nucleus, thalamus, mesial inf. temporal lobe
\item occipital and occipitoparietal cortex
\item Affect vision
\item contralateral homonymous hemianopsia
\item cortical blindness
\item Visual agnosia
\item Altered mental status, Impaired memory
\end{itemize}
\end{latin}
\IPfoot{۱۹}
% pic20.jpg
\IPpage{۲۰}
\begin{latin}
\noindent Vertebrobasilar system
\begin{itemize}
\item medulla, cerebellum, pons, midbrain, thalamus, occipital cortex
\item Major disability or death
\item Crossed findings: Ipsilateral cranial nerve deficits + Contralateral motor deficits
\end{itemize}
\end{latin}
\subsection*{سکته هموراژیک}
\begin{enumerate}[label=\arabic*.,leftmargin=*]
\item ۱۵\% موارد سکته‌ها هموراژیک است.
\item پروگنوز آن از ایسکمیک بدتر است.
\item خونریزی می‌تواند در پارانشیم و یا در فضای ساب‌آراکنوئید رخ دهد.
\item اکثر بیماران \en{HTN} دارند.
\item ۹۰\% موارد خونریزی‌های ناشی از \en{HTN} در بازال گانگلیا رخ می‌دهد.
\begin{IPpoints}\item دیگر نواحی: تالاموس، \en{Pons}، مخچه\end{IPpoints}
\item درمان باید ظرف \en{4--5 h} از شروع علائم آغاز گردد.
\end{enumerate}
\IPfoot{۲۰}
% pic21.jpg
\IPpage{۲۱}
\noindent ۷. تصمیم اینکه سکته ترومبولیتیک یا ریکانالیزیشن بکنیم وابسته به یافته‌های \en{NC-CT} اولیه است.

\noindent ۸. \en{Recognizing ICH}
\begin{IPenglish}
\item Freshly extravasated whole blood: Increased density on NC-CT; Due to hemoglobin
\item Dissection of blood into the ventricular system: In HTN bleeds
\item As the clot begins to form, the blood becomes denser for about 3 days; Due to dehydration of the clot
\item After the third day: the clot decreases in density; Invisible over the next several weeks
\item After about 2 months: Only a small hypodensity may remain
\end{IPenglish}
\subsection*{آنوریسم پاره‌شده}
\noindent ۱. شایع‌ترین آنوریسم در سیستم عصبی: آنوریسم \en{Berry}
\begin{IPpoints}
\item به علت ضعف مادرزادی در دیواره شریان
\item شاخه‌های حلقه ویلیس را درگیر می‌کند.
\item می‌تواند فامیلیال یا مرتبط با بیماری‌های بافت همبند باشد.
\end{IPpoints}
\noindent ۲. سیگار و \en{HTN}: ریسک آنوریسم‌ها
\IPfoot{۲۱}
% pic22.jpg
\IPpage{۲۲}
\noindent ۳. آنوریسم هر چه بزرگ‌تر باشد احتمال خونریزی آن بیشتر است.

\noindent ۴. \en{Lacunar Infarcts}
\begin{IPpoints}
\item انفارکت‌های مغزی کوچک
\item ناشی از انسداد شریان‌های کوچک انتهایی
\item مسئول ۲۰\% موارد انفارکت مغزی
\item شایع‌تر در: بازال گانگلیا، اینترنال کپسول، \en{Pons}
\item در همراهی با \en{HTN} رخ می‌دهد؛ دیابت و آترواسکلروز
\end{IPpoints}
\noindent ۵. \en{Chronic Lacunar Infarct}
\begin{IPenglish}\item low-density, cystic lesions about 5 to 15 mm in size.\end{IPenglish}
\noindent ۶. \en{Appearance of blood over time on MRI}
\begin{latin}
\begin{center}
\begin{tabular}{llll}
Stage & Time & $T_1$ & $T_2$ \\
\hline
Hyperacute & $<24$ hours & Isointense & Bright \\
Acute & 1--3 days & Isointense & Dark \\
Early subacute & 3--7 days & Bright & Dark \\
Late subacute & 7--14 days & Bright & Bright \\
Chronic & $>14$ days & Dark & Dark \\
\end{tabular}
\end{center}
\end{latin}
\IPfoot{۲۲}
% pic23.jpg
\IPpage{۲۳}
\noindent ۷. اندازه حیاتی آنوریسم \en{10 mm} است.

\noindent ۸. شرح حال کلاسیک بیمار آنوریسم: بدترین سردرد عمرم!

\noindent ۹. در پی پاره شدن آنوریسم خون وارد فضای ساب‌آراکنوئید می‌شود.

\noindent ۱۰. پارگی آنوریسم شایع‌ترین علت غیرترومایی \en{SAH} است. (۸۰\%)

\noindent ۱۱. آنوریسم‌ها با \en{CTA} و \en{MRA} تشخیص داده می‌شوند.
\begin{IPpoints}\item با کنتراست / بدون کنتراست\end{IPpoints}
\noindent ۱۲. \en{Recognizing SAH from a rupture of aneurysm}
\begin{IPenglish}
\item Acute blood in CT: Hyperdense; In Sulci and Basal cisterns
\item Region of the falx: Hyperdense, Widened, Irregularly marginated
\end{IPenglish}
\noindent ۱۳. به طور کلی در آنوریسم محل تجمع خون نشان‌دهنده محتمل‌ترین محل آنوریسم است.
\subsection*{هیدروسفالی}
\noindent ۱. علل
\begin{IPpoints}
\item تولید بیش از حد مایع \en{CSF} (ناشایع)
\item کاهش جذب \en{CSF} (هیدروسفالی کامیونیکیتینگ)
\item محدودیت خروج \en{CSF} از بطن‌ها (هیدروسفالی نان‌کامیونیکیتینگ)
\end{IPpoints}
\IPfoot{۲۳}
% pic24.jpg
\IPpage{۲۴}
\noindent ۲. افتراق بزرگ شدن بطن‌ها و سولکوس‌ها در هیدروسفالی و آتروفی مغز
\begin{IPpoints}\item در هیدروسفالی بطن‌ها نسبت به سولکوس‌ها بیشتر بزرگ می‌شوند.\end{IPpoints}
\noindent ۳. در هیدروسفالی اندازه شاخ‌های تمپورال (\en{Temp. horns}) بیشتر از \en{2 mm} می‌شوند و در تصویربرداری‌ها قابل رؤیت می‌گردند.

\noindent ۴. هیدروسفالی انسدادی
\begin{IPenglish}\item Communicating (extraventricular obstruction)\item Noncommunicating (intraventricular obstruction)\end{IPenglish}
\noindent ۵. هیدروسفالی ارتباطی
\begin{IPpoints}\item اختلال در سطح ویلوس‌های آراکنوئید\item بطن چهارم دیلاته است.\end{IPpoints}
\noindent ۶. هیدروسفالی غیرارتباطی: بطن چهارم دیلاته نیست.

\noindent ۷. درمان هیدروسفالی ارتباطی: \en{Ventricular Shunt}

\noindent ۸. علل هیدروسفالی غیرارتباطی
\begin{IPpoints}\item تومورها\item کیست‌ها\item سایر موانع فیزیکی\end{IPpoints}
\noindent ۹. هیدروسفالی مادرزادی: انسداد بین بطن‌های سوم و چهارم در سطح آکواداکت سیلویوس
\IPfoot{۲۴}
% pic25.jpg
\IPpage{۲۵}
\noindent ۱۰. اگر علت انسداد تومور یا کیست باشد: جراحی

\noindent ۱۱. هیدروسفالی غیرانسدادی ناشی از تولید بیش از حد \en{CSF}
\begin{IPpoints}\item بسیار نادر\item کوروئید پلکسوس پاپیلوما\end{IPpoints}
\noindent ۱۲. هیدروسفالی با فشار طبیعی: \en{Normal Pressure Hydrocephalus (NPH)}
\begin{IPpoints}
\item فرمی از هیدروسفالی ارتباطی است.
\item تریاد تشخیصی: آبنورمالی گیت، دمانس، بی‌اختیاری ادرار
\item سن شروع معمولاً بین ۴۰ تا ۷۰ سال است.
\item تشخیص آن مهم است چون درمان خوبی دارد: شانت ونتریکولوپریتونئال
\item یافته‌ها در تصویربرداری: بزرگ شدن بطن‌ها؛ بزرگ شدن تمپورال هورن‌ها؛ سولکوس‌ها نرمال یا تخت
\end{IPpoints}
\subsection*{آتروفی مغز}
\noindent ۱. آتروفی مغز: از دست رفتن ماده سفید و خاکستری می‌شود.

\noindent ۲. یافته اصلی ولی غیراختصاصی در آلزایمر
\begin{IPpoints}\item آتروفی کورتیکال مغز\item هیپوکامپ تمپورال\end{IPpoints}
\noindent ۳. در پی آتروفی، سولکوس‌ها و بطن‌ها بزرگ می‌شوند.
\IPfoot{۲۵}
% pic26.jpg
\IPpage{۲۶}
\subsection*{تومورهای مغز}
\noindent ۱. گلیوم مغز
\begin{IPpoints}
\item شایع هستند: ۳۰\% همه تومورهای مغز؛ ۸۰\% همه بدخیمی‌های اولیه مغز
\item \en{primary supratentorial intraaxial mass}
\item نیمی از موارد گلیوم‌ها: گلیوبلاستوم مولتی‌فرم؛ لوب فرونتال و تمپورال
\item ۲۰\% از موارد گلیوم: آستروسایتوما؛ مردان ۲۵ تا ۴۵ سال
\item گلیوبلاستوم مولتی‌فرم: بدترین پروگنوز بین گلیوم‌ها
\item به اطراف دست‌اندازی می‌کند.
\item متاستاز خارج مغزی نمی‌دهد.
\end{IPpoints}
\noindent ۲. تشخیص گلیوبلاستوم مولتی‌فرم
\begin{IPpoints}
\item چون رشد تهاجمی دارد: نکروز در وسط تومور
\item بافت مغز را انفیلتراسیون می‌دهد و از ماده سفید عبور می‌کند.
\item از کورپوس کالوزوم عبور به نیمکره مقابل: \en{Pattern of Butterfly glioma}
\item ادم وازوژنیک و اثر فشاری ایجاد می‌کند.
\item با ماده کنتراست \en{enhance} می‌شود.
\end{IPpoints}
\noindent ۱. شایع‌ترین تومورهای متاستازدهنده به مغز: ریه، پستان، ملانوم
\IPfoot{۲۶}
% pic27.jpg
\IPpage{۲۷}
\noindent ۲. \en{Recognizing metastasis to the brain}
\begin{IPenglish}
\item Well defined round masses near the gray-white junction.
\item Usually multiple
\item Typically hypodense or isodense on NC-CT
\item With contrast: Ring enhancement pattern
\item Vasogenic edema
\end{IPenglish}
\noindent ۳. مننژیوم
\begin{IPpoints}
\item شایع‌ترین تومور اولیه مغز
\item شایع‌ترین توده اکسترا‌آگزیال
\item معمولاً در زنان میان‌سال
\item مکان‌های شایع: پاراساژیتال؛ بال اسفنوئید؛ \en{Over the convexities}؛ \en{CPA} (زاویه سربلوپونتین)
\item رشد آهسته دارند.
\item درمان جراحی؛ پروگنوز خیلی عالی
\item اگر متعدد باشند ممکن است با نوروفیبروماتوز تیپ \en{II} مرتبط باشد.
\end{IPpoints}
\noindent ۴. \en{Recognizing meningioma on CT scan}
\begin{IPenglish}
\item $>50\%$: hyperdense on NC-CT; $20\%$ contain calcification
\item On enhanced CT they enhance markedly.
\item Vasogenic Edema
\end{IPenglish}
\IPfoot{۲۷}
% pic28.jpg
\IPpage{۲۸}
\noindent ۴. شوانوم وستیبولار (آکوستیک نوروما)
\begin{IPpoints}
\item شایع‌ترین شوانوم اعصاب کرانیال
\item علائم: ناشنوایی، صدای وزوز گوش، \en{Disturbance in equilibrium}
\item معمولاً درون \en{Internal auditory canal} و \en{CPA}
\item اگر متعدد باشند ممکن است با \en{NF2} مرتبط باشند.
\end{IPpoints}
\subsection*{مالتیپل اسکلروزیس}
\noindent ۱. شایع‌ترین بیماری دمیلینیتینگ: \en{MS}

\noindent ۲. هر عملکرد نورولوژیکی می‌تواند تحت تأثیر \en{MS} قرار بگیرد.

\noindent ۳. دوره‌های عودکننده و خاموش‌شونده دارد.

\noindent ۴. مسیرهای میلینه در ماده سفید را درگیر می‌کند که به ضایعات آن پلاک می‌گویند.

\noindent ۵. نواحی ضایعات
\begin{IPpoints}\item ناحیه پری‌ونتریکولار\item کورپوس کالوزوم\item عصب اپتیک\end{IPpoints}
\noindent ۶. \en{MRI} تصویربرداری انتخابی است.
\begin{IPenglish}
\item Discrete globular foci of high-signal intensity on $T_2$
\item In acute MS, lesions enhance with gad on $T_1$.
\item Oriented with their long axes perpendicular to the ventricular walls
\end{IPenglish}
\IPfoot{۲۸}
\ifnum\IPStandalone=1 \end{document}\fi
